\subsection{Amplitudes over arbitrary supports with semiperfect upper rings}
We use the ordinary semiperfect and finite Krull--Schmidt theorems
in \cite[Proposition~4.1, pp.~8--9, Theorem~4.2 and Corollary~4.3,
p.~9 of the author preprint]{Krause}. Their application below is to
a finite coefficient ring constructed from the projector, not to
the entire action category.

\begin{proof}[Proof of Theorem~\ref{thm:upper-semiperfect-amplitudes}]
Apply Theorem~\ref{thm:finite-coefficient-models} to $G$. It gives
an actual unital coefficient ring $C=A_D(T)$, an ideal $N\lhd C$,
a map $\pi:C\to B$ to a finite full action-arrow ring, and an
idempotent $e\in N$ representing $G$. Every mixed corner between
idempotents in $N$ maps isomorphically to its full action corner.
The diagonal corners of $C$ are $F(d_{ii})$ for upper idempotents
$d_{ii}$. Decompose each of their identities into finitely many
orthogonal idempotents with local corners. Together these form
such a decomposition of $1_C$, so $C$ is semiperfect. This step
makes no assertion about semiperfectness of $B$.

Put $I=CeC\subseteq N$. The finite projective left $C$-module
$Ce$ is $I$-generated. Ordinary semiperfect theory decomposes it
into columns $Cf$ from the chosen diagonal local-idempotent list.
Ideal generation passes to each occurring summand, so
$Cf=ICf\subseteq I$ and in particular $f\in I\subseteq N$.
The exact-corner assertion gives
\[
 fCf\simeq\pi(f)B\pi(f).
\]
Thus $\pi(f)$ is nonzero and has a local full action corner.
It lies in a single diagonal object and in the canonical ideal.
The decomposition of $Ce$ supplies rectangular idempotent
intertwiners over $C$. Applying $\pi$ and composing on
representables identifies $G$ with the corresponding sum of
$\mathcal A(l_i,-)\pi(f)$. These inverse matrices require no
flatness assumption. This proves existence of the asserted
summands using only an ordinary coefficient ring and exact
corners.

Conversely $\mathcal A(l,-)f$ is the image of the idempotent $f$
on one representable. It is $\mathcal J$-generated since its
generator $f$ belongs to $\mathcal J(l,l)$ and $f=f^2$.
Its endomorphism ring is the opposite of the nonzero local ring
$f\mathcal A(l,l)f$. The Yoneda corner calculation identifies its
isomorphisms exactly with the two displayed inverse-arrow
relations. Krull--Schmidt uniqueness for finite sums of objects
with local endomorphism rings proves uniqueness of the
multiplicities. Theorem~\ref{thm:canonical-ideal} now gives the
semimodule classification, and transport of $\mathcal A$ and
$\mathcal J$ gives precisely the support-automorphism orbits.

For the local specialization, let $0\ne f=f^2\in\mathcal J(l,l)$.
Represent $f$ at $q=\phi(l)$ for an admissible $\phi$. As above,
$q^2\leq q$ and $e=1\vee q$ is an upper idempotent stabilizing
$l$. The unital image of $F(e)$ in $\mathcal A(l,l)$ is local
and contains $f$. A nonzero idempotent in a local ring is its
identity, so $f=1_l$. Conversely an allowed identity is a
permitted local idempotent, as the next paragraph verifies.
\end{proof}

\begin{corollary}[the local diagonal alternative]
\label{cor:upper-local-diagonal}
Suppose all upper coefficient rings are local or zero. Put
$A_l=\mathcal A_L(l,l)$ and $I_l=\mathcal J_L(l,l)$.
If $I_l=A_l$, then $A_l$ is local or zero. Otherwise
$I_l\subseteq\operatorname{Jac}(A_l)$.
\end{corollary}
\begin{proof}
An element of $I_l$ is represented at $q=\phi(l)$ for an
admissible $\phi$, and hence lies in the unital local image of
$F(1\vee q)$ in $A_l$. If $I_l$ is proper, such an element
cannot be a unit in that image. Its complement to $1$ is therefore
a unit in $A_l$. Applying this to every element of the ideal
$I_l$ gives the ordinary Jacobson criterion.

If $1_l=[q,x]\in I_l$, represent any $a\in A_l$ by $[p,y]$
with $pl\leq l$. Then $pq=\phi(pl)\leq q$, and
$a=a1_l=[pq,yx]$ transports to $F(q)$. Thus the unital image of
$F(1\vee q)$ is all of $A_l$, proving the assertion.
\end{proof}

\begin{example}[an action ring need not be local]
\label{ex:upper-local-action-boundary}
Let $Q=\{\bot\}\sqcup\mathbb N$, with maximum for join,
$m n=m+n$ on $\mathbb N$, absorbing $\bot$, and unit the integer
$0$. For a field $k$, put $F(\bot)=0$ and let $F(n)$ be the
additive group of polynomials in $k[t]$ of degree at most $n$.
Transitions are inclusions and products are ordinary polynomial
products. These satisfy all pair axioms, with unit $1\in F(0)$.
The sole upper multiplicative idempotent is the integer $0$, whose
ring is $k$.

For the singleton zero $Q$-support, the action endomorphism ring is
$\varinjlim_nF(n)=k[t]$, which is not local: $(t)$ and $(t-1)$ are
distinct maximal ideals. The only admissible functional is the
zero functional with value $\bot$, so the canonical ideal is zero
and only zero amplitude is permitted. The domain $k[t]$ has only the idempotents $0,1$, so its
semiperfectness would imply locality. Thus local upper fibers do
not make all action endomorphism rings local or semiperfect.
\end{example}

